\documentclass{nfr}

% This LaTeX template is based on the official Word template found here: https://www.forskningsradet.no/siteassets/utlysninger/vedlegg-utlysninger/template-for-project-description-researcher-project-for-experienced-scientists-and-early-career-scientists.docx
% See the README.md file for more information

% Here's an example bibliography using the numeric (Nature) style and two
% columns for the reference list.

% \usepackage{multicol}
% \renewcommand{\bibpreamble}{\begin{multicols*}{2}}
% \renewcommand{\bibpostamble}{\end{multicols*}}

% \usepackage{biblatex}
% \addbibresource{references1.bib}

% \citestyle{nature}
% \citestyle{}
\usepackage{comment}

% \bibliographystyle{bst/veryshortnum}
\bibliographystyle{apalike}

\title{Application for Postdoctoral Research Position in Marketing at the University of Basel}
\author{Zhongli Wang}

\begin{document}

% Project description template
% Researcher Project for Experienced Scientists/Early Career Scientists

% This template consists of
% 
% \textbf{A:} Overview of the chapters and sections to be used as the
% structure of the project description
% 
% \textbf{B:} General guidance
% 
% \textbf{C:} Guidance on the content in each chapter and section
% 
% \textbf{A: Chapters and sections to be used}

% I used pandoc to translate the word template to tex. The formatting of the
% appendices are, therefore, not great, but the contents are there.

\maketitle

% Understanding firm competition in modern markets remains a fundamental challenge in economics. 
% While classical models such as Bertrand, Cournot, and Hotelling provide valuable benchmarks, real markets are characterized by firm asymmetries, network effects, and heterogeneous consumers, which complicate predictions and strategic decision-making \citep{Hotelling1929}.
% My research tackles this challenge by combining computational and statistical modeling to study strategic behavior and market dynamics in complex economic environments. Using tools from agent-based modeling, industrial organization, and econometrics, I analyze how firms and consumers interact in digital and networked markets. The broader goal is to bridge theory, simulation, and data to generate solutions that improve market understanding and support effective policy and managerial decisions.

\noindent Dear Dr.~Bublitz and Members of the Hiring Committee,\\[0.2cm]

I am writing to apply for the Postdoctoral Researcher position in Marketing (AI \& Synthetic Data) at the University of Basel. 
I am currently completing my Ph.D.\ in Economics at Bielefeld University and Università Cattolica del Sacro Cuore (EPOC Program, Marie Skłodowska-Curie Grant No.~956107). 
My research lies at the intersection of computational social science, machine learning, and applied microeconomics, with a strong focus on designing AI-assisted pipelines for data generation, behavioral modeling, and policy evaluation. 
I am particularly interested in how large-language-models (LLMs) can produce reliable synthetic data for understanding consumer behavior, communication dynamics, and opinion formation—questions that align closely with the objectives of your Innosuisse project and with your research on persuasion and misinformation.\\

I began my Ph.D.\ in 2021, during a period in which the COVID-19 pandemic exposed how quickly digital infrastructures and information environments can change. 
This context shaped my dissertation, which develops computational frameworks to analyze how firms and consumers behave in technology-intensive markets. 
In my central dissertation chapter, I build and estimate an agent-based model of ICT platforms in which heterogeneous users respond to service quality, congestion, and information flow. 
Using firm-level data, simulation-based econometrics, and machine-learning tools, I estimate key behavioral parameters and evaluate counterfactual technology-adoption scenarios. 
Across all three chapters—each single-authored—I integrate microeconomic modeling with modern computational methods to generate, validate, and analyze synthetic market data. 
These techniques translate naturally to LLM-based synthetic data generation and validation, where questions of behavioral realism, data fidelity, and model calibration are equally central.\\

Beyond my dissertation, my recent coauthored papers—\textit{Polytope Fraud Theory} and \textit{Agentic Workflows for Economic Research}—extend my methodological work toward computational social science and AI-assisted data generation \citep{Zhao2025, Dawid2025}. 
The first paper employs machine-learning techniques to detect behavioral anomalies in large-scale financial datasets, demonstrating how statistical learning can reveal latent patterns and support data-driven decision processes. 
The second paper develops a general framework for building LLM-based agentic workflows, including retrieval-augmented generation, prompt engineering, and model evaluation. 
This work shows how LLM pipelines can be systematically designed, audited, and validated—an approach that directly informs ongoing efforts to generate reliable synthetic data for marketing and opinion research.\\

My primary training is in economics, yet my research employs methods that are highly relevant to marketing analytics, including large-scale data handling, machine-learning pipelines, and behavioral modeling. 
I have also developed a strong interest in marketing and consumer research. 
In 2024, I participated in the “Challenges in Managing Smart Products and Services (CHIMSPAS)” conference hosted by the Institute of Technological Innovation, Market Development and Entrepreneurship (iTIME) at Bielefeld University, where I presented the first chapter of my dissertation and received valuable feedback from scholars working on digital products, consumer behavior, and technology adoption.
In my future work, I aim to build on recent advances such as \citet{arora2025ai}, which demonstrate that LLMs can assist in both data generation and data analysis by creating realistic sample characteristics, producing synthetic respondents, and even conducting or moderating in-depth interviews. 
I am particularly interested in extending these ideas toward more rigorous validation frameworks for marketing and opinion research. 
At present, I am focusing on improving retrieval-augmented generation (RAG) pipelines and integrating human-in-the-loop procedures to ensure that LLM-generated responses exhibit behavioral realism, internal consistency, and empirical validity.\\

In parallel with my dissertation work, I have gained substantial experience managing large datasets and coordinating multi-institutional collaborations. 
My doctoral research, funded by the \textit{Marie Skłodowska-Curie Actions} and the \textit{Bielefeld Young Researchers’ Fund}, involved designing and executing complex empirical and computational workflows, integrating simulation results, machine-learning models, and external data sources. 
Through these projects, I developed strong skills in project coordination, reproducible data management, and collaboration with academic and industry partners—competencies that I would bring to the Innosuisse project and to joint work with Boomerang Ideas, SBB, and gfs.bern.\\

Looking ahead, my research agenda focuses on developing rigorous and scalable methods for LLM-based synthetic data generation, with applications to marketing analytics, communication research, and opinion formation. 
I am particularly interested in advancing validation frameworks that combine survey experiments, retrieval-augmented generation, and human-in-the-loop procedures to assess the behavioral realism and reliability of LLM-generated responses. 
In the medium term, I hope to contribute to competitive funding initiatives in Switzerland and Europe by developing projects that integrate synthetic data, causal inference, and computational social science to address challenges in persuasion, misinformation, and digital communication. 
I am excited about the possibility of pursuing these ideas within Basel’s strong research environment and in collaboration with the Innosuisse project partners.\\

Thank you so much for considering my application! 
I would be honored to join the University of Basel as a postdoctoral researcher and to contribute to your work on LLM-assisted synthetic data generation in marketing and opinion research. 
I am especially excited about the opportunity to collaborate with you and the Innosuisse project partners, and to develop rigorous, scalable methods that advance both academic research and practical applications in industry. 
I greatly look forward to the possibility of contributing to Basel’s dynamic research environment.
\vspace{3mm}

\noindent Sincerely,\\
\vspace{10mm}

\noindent Zhongli Wang

\vspace{5mm}

% Throughout my doctoral and collaborative research, I have experienced the value of combining analytical modeling with computational innovation. Working across fields has shown me that the most pressing problems—such as market transparency, sustainable technology adoption, and responsible data governance—can only be solved through interdisciplinary collaboration.

% In addition to my market-oriented research, my coauthored publication \textit{``Polytope Fraud Theory''} and the working paper \textit{``Agentic Workflows for Economic Research''} have shaped my commitment to methodological innovation. The former integrates economic reasoning with machine learning to enhance fraud detection and investor protection; the latter connects computational economics with artificial intelligence by designing LLM-based research workflows. These experiences have strengthened my belief that economists can contribute not only to theoretical insights but also to the design of intelligent, transparent, and reproducible systems.

% I am therefore particularly excited about joining the FINATRAX Research Group, where business research meets information systems engineering. The group’s mission—to explore digital identities, distributed ledgers, machine learning, and IoT as enablers of sustainable technologies—closely aligns with my research vision. I am enthusiastic about participating in the design, implementation, and theoretical elaboration of high-quality, interdisciplinary solutions that create tangible value for both organizations and society. This position represents a natural continuation of my efforts to bridge economics with emerging digital technologies.

% \section*{Funding and Collaboration Strategy}

% My strategy for research funding and partnerships builds on a dual foundation of analytical rigor and applied collaboration. 
% Methodologically, I plan to combine analytical and computational approaches: analytical derivations will establish benchmark propositions, while numerical simulations and agent-based modeling will explore market dynamics and firm behavior under varying degrees of information transparency. This structure naturally supports both theoretical innovation and data-driven applications—key elements for competitive research proposals.

% For empirical validation, I will leverage access to the Refinitiv database through collaboration with the Southern University of Science and Technology (SUSTech), complemented by additional data resources available at the host institution. My doctoral research, funded by the \textit{Marie Skłodowska-Curie Actions} and the \textit{Bielefeld Young Researchers’ Fund}, has already demonstrated my ability to manage funded projects and coordinate cross-institutional collaboration. 

% Going forward, I intend to apply for European and national-level grants such as \textit{Horizon Europe}, \textit{DFG}, and \textit{VolkswagenStiftung}, focusing on projects that integrate computational economics with digital transformation research. On the industry side, I aim to build partnerships with firms in the ICT and FinTech sectors to co-develop decision-support tools and simulation frameworks for market design and policy evaluation. These collaborations would ensure practical relevance, foster technology transfer, and open opportunities for joint funding applications. With adequate computational resources and institutional support, I am confident that these efforts can yield both high-impact publications and meaningful contributions to industry and policy communities.





\begin{comment}

\section*{Statement of Research Interests}

% Here's a citation to \cite{texbook}. Here are citations to \cite{texbook,knuth:1984,lesk:1977}, and to \cite{lesk:1977,texbook,latex2e,latex:companion}.

My doctoral dissertation, ``\textit{Strategic Behavior and Economic Dynamics in ICT Markets: An Agent-based Framework}", investigates how firms in the information and communications technology (ICT) industry adapt under constraints such as data regulations, capacity limits, and consumer network effects. The work builds a bridge between state-space and agent-based models, showing how techniques from dynamic latent-variable systems can be applied to simulate and estimate complex economic environments.\footnote{All three chapters are single-authored papers.}\\

My job market paper, \textit{``Strategic Firm Capacity Selection under Network Effects: A Computational Model of Capacity Investment Dynamics in the ICT Industry''}, investigates how digital firms determine optimal service capacity under the joint influence of consumer congestion and network effects---an interaction that has received limited attention in the existing literature. While traditional industrial organization models emphasize innovation or price competition, many digital firms face a different constraint: their profitability depends critically on infrastructure investment and the ability to maintain service quality. To address this challenge, I combine game-theoretic reasoning with agent-based simulations. The game-theoretic framework captures the multi-stage nature of firms’ decisions on pricing and capacity investment, while the agent-based model represents the market dynamics that emerge when heterogeneous consumers interact through a small-world network and share information within user communities. The results show that strong network effects tend to amplify competitive advantages by concentrating demand, whereas congestion generates counteracting pressures that reduce pricing flexibility and erode profits. Together, these forces explain why digital firms must balance capacity expansion and user experience to sustain long-term profitability. This study is among the first to systematically analyze the joint role of network and congestion effects in the ICT industry, offering new insights into pricing, investment, and regulatory strategies for digital service markets.\\

Building on the similar market context, my previous paper, \textit{``Heterogeneous Consumers, Firm Competition, and the Data Value Chain: An Agent-Based Approach''}, investigates how differences in data-governance intensity influence competition in software markets. The challenge arises because digital firms depend on user data to enhance product quality and pricing decisions, while local governments impose varying levels of data regulation to protect privacy and maintain market fairness. 
To analyze this tension, I develop a hybrid game-theoretic and agent-based model in which governments determine the stringency of data governance, and firms subsequently compete on investment and pricing decisions. 
Consumers are heterogeneous in their sensitivity to product quality and their attitudes toward data governance, which shape their purchasing behavior. 
The results show that stricter data governance limits firms’ ability to utilize consumer data effectively and reduces competitive differentiation, which leads to higher prices and lower consumer welfare. In contrast, moderate levels of regulation can maintain competition and promote more balanced market outcomes. The analysis suggests that regional governments would benefit from greater coordination in setting data policies and from incorporating consumer preferences into policy design, thereby supporting both market efficiency and innovation. Overall, these findings offer a structured perspective on how policymakers can balance privacy protection with incentives for innovation in digital markets.\\

Building on the insights from the first two studies, my third paper ``\textit{Parameter Estimation Methods in an Agent-based Computational Model: Insights on ICT Industrial Dynamics and Policies}" focuses on empirical validation. 
Using financial and operational data from two leading U.S. telecommunication firms, AT\&T and T-Mobile, I estimate the model parameters with the Simulated Method of Moments and neural network–based surrogate modeling (both within-sample and out-of-sample check). 
The main challenge lies in linking complex agent-based models to real-world data, as repeated simulations are computationally intensive and traditional estimation methods often fall short. 
To address this, I employed the Paderborn High-Performance Computing (HPC) clusters and the Python package \texttt{ray} to parallelize the simulations and collect 40{,}000 samples for parameter search. 
By comparing these two approaches, the study provides practical evidence on how to calibrate computational models efficiently while maintaining empirical accuracy. This work contributes to the broader methodological discussion on integrating simulation-based models with data-driven analysis to better understand market behavior and policy outcomes.



\section*{Publications and Emerging Projects}

In addition to my market-oriented research, I have one peer-reviewed publication and one current working paper that contribute to the broader use of computational methods in economics.\footnote{Both papers are coauthored with other researchers.} \\

The paper \textit{``Polytope Fraud Theory''} extends the classical Fraud Triangle and Fraud Diamond frameworks by introducing a multidimensional structure that captures ten abnormal financial practice alarms indicative of potential corporate fraud \citep{Zhao2025}. Using data on activist short sellers and supervised machine learning techniques, the study identifies financial patterns associated with fraudulent reporting and demonstrates how algorithmic approaches can support market transparency and investor protection. 
The paper also demonstrates the potential for automating advanced financial statement analysis, providing tools that can enhance both auditing efficiency and investment decision-making. In addition, it introduces the \textit{Unified Investor Protection Framework}, which systematically categorizes investor-protection theories from the macro- to the micro-level. This framework offers practical guidance for identifying firm-level financial risks and strengthens the connection between analytical modeling and real-world investor decision processes.\\

Another working paper, \textit{``Agentic Workflows for Economic Research: Design and Implementation''}, develops a framework that employs large-language-model (LLM) agents to support economists throughout the research process—from data collection and coding to simulation design and interpretation \citep{Dawid2025}. The framework formalizes how autonomous AI agents can collaborate within multi-stage research pipelines while ensuring reproducibility, version control, and transparent documentation. Using examples implemented in Microsoft AutoGen, we demonstrate how such workflows can accelerate tasks like research-question ideation, literature review, and model specification, thereby improving research efficiency without compromising academic rigor. Importantly, the framework positions human researchers at the core of the process: AI agents act as supportive tools that augment, rather than replace, human judgment and creativity. This line of work extends my methodological interests in AI-enhanced modeling, bridging computational economics with artificial-intelligence research and tool development.

\section*{Future Research}

\subsection*{Potential Research Direction: The Observability of Firm Heterogeneity in Price-Directed Consumer Search}

% Recent models of price-directed search \citep{obradovits2023price} highlight how consumers allocate their attention and search effort across differentiated firms when facing information frictions. However, these frameworks often assume that firm heterogeneity---in cost structures, quality, or service intensity---is either fully observable to consumers or symmetrically distributed across sellers. In many real markets, such as digital platforms, online retailing, or app ecosystems, this heterogeneity is only partially observable: consumers infer firm characteristics indirectly from prices, reputation indicators, or platform signals.

% Understanding how the observability of firm heterogeneity shapes consumer search behavior and market efficiency is crucial for both theory and policy. It determines the extent of equilibrium price dispersion, the efficiency of market matching, and firms’ incentives to disclose information. The problem is empirically feasible because modern datasets (e.g., online pricing panels, platform APIs, or laboratory experiments) provide micro-level evidence on search behavior under varying degrees of information transparency. Computational approaches, including agent-based simulations, further enable the exploration of equilibrium dynamics in complex environments.

Recent models of price-directed search \citep{obradovits2023price} examine how consumers allocate attention and search effort across differentiated firms under information frictions, yet they typically assume that firm heterogeneity---in costs, quality, or service intensity---is fully observable or symmetrically distributed. In real markets such as digital platforms or online retailing, this heterogeneity is only partially observable, as consumers infer firm characteristics indirectly from prices, ratings, or platform signals. Understanding how such observability affects consumer search and market efficiency is crucial, as it shapes equilibrium price dispersion, allocative outcomes, and firms’ incentives to disclose information. The question is empirically feasible given modern micro-level data on online search and pricing, while computational approaches, including agent-based simulations, allow systematic analysis of equilibrium dynamics under varying information transparency. \\

% This research direction contributes to Industrial Organization by linking information structure, firm heterogeneity, and price competition under bounded rationality. It combines theoretical modeling with simulation and empirical validation, making it both conceptually rich and methodologically tractable. 

The following hypotheses will be tested:

\begin{enumerate}%[label=\textbf{H\arabic*.}, leftmargin=2em]
    \item \textbf{Partial observability increases price dispersion.} When firm heterogeneity (e.g., in cost or quality) is imperfectly observable, consumers engage in less directed search, resulting in a wider equilibrium price distribution than under full observability.
    
    \item \textbf{High observability improves allocative efficiency.} Markets with more transparent signals of firm heterogeneity (such as ratings or quality certifications) achieve a closer match between consumers and preferred firms, raising consumer surplus and reducing welfare loss.
    
    \item \textbf{Strategic transparency depends on firm type.} Low-cost or high-quality firms have stronger incentives to increase the observability of their characteristics (e.g., through signaling or advertising) than high-cost or low-quality competitors, leading to asymmetric information equilibria.
\end{enumerate}

\subsection*{Methodological Plan}

In my approach, I plan to combine analytical and computational methods. Analytical derivations will be used to establish the basic propositions and benchmark outcomes of the model, while numerical techniques and computational modeling will be employed to solve for Nash equilibria and explore market dynamics under varying degrees of information transparency. For empirical validation, I have access to the Refinitiv database through an application via the Southern University of Science and Technology (SUSTech) and may also rely on additional database resources at my future host institute. My doctoral research was supported by the Marie Skłodowska-Curie Actions and the Bielefeld Nachwuchsfonds, and it benefited from extensive computational work conducted on Google Colab (self-funded) and the Paderborn High Performance Computing cluster (supported by Bielefeld University). I hope that the new institute could provide adequate support for research infrastructure, particularly computational resources, to further advance this line of study.



% \subsection*{Computational \& ABM Dynamics}

% Specifically, I plan to:

% •	Develop scalable frameworks for data-driven research of economic models that can inform competition and innovation policy.

% •	Study how AI and automation transform firm strategy, market structure, and industrial dynamics.

% My long-term objective is to contribute to the economics literature, that is both analytically rigorous and computationally adaptive—one that captures the rich, evolving behavior of firms and consumers in real markets. I remain deeply motivated to pursue research that connects theoretical innovation with practical relevance, and I look forward to collaborating in an environment that values interdisciplinary approaches and empirical credibility.
% \citep{arve2025decentralized}.


% \bibliography{NFR_Template_Statements/references}

\newpage

\appendix
\end{comment}
\bibliography{NFR_Template_Statements/references1}

\end{document}
